\section{Introduction}\label{sec:intro}
This paper concerns the  problem of determining whether a formula in some logic $\mathcal{L}$ is satisfied in some finite structure: the \emph{finite satisfiability problem for} $\mathcal{L}$. For the fragments of interest to us, the decidability of this problem will imply decidability of many other static analysis problems for the logic, such as validity of a sentence over finite structures and equivalence of two sentences over finite structures.
The problem is known
to be undecidable for first order logic~\cite{traktenbrot}. 
In the last decades two paradigms for decidability emerged, one based on restricting to \emph{guarded quantification}, and another restricting to \emph{two variable fragments}.
The two variable fragment of first order logic was shown to have decidable finite satisfiability problem and this was extended to allow ``counting quantifiers'',  of the form $\exists^{\ge k} y ~ \phi(x,y)$, where $k$ is fixed. The complexity of the satisfiability problem for this logic, normally denoted $\Ctwo$, over finite structures, as well as the complexity of satisfiability when considering arbitrary structures, was isolated in~\cite{ctwocomplexity}. The same problem for the sublogic where quantification is guarded was shown 
to have slightly lower complexity~\cite{gctwocomplexity}.

In recent years there has been interest in adding more powerful arithmetic capabilities to these decidable logics. Instead of just asking whether the number of elements satisfying a property is above or below a constant, 
we can ask whether the cardinality is even, or whether 
the cardinality of several properties satisfies a linear inequality. 
\emph{Presburger logic} focuses on linear inequalities between properties, 
and \emph{local Presburger logic}, $\GPtwo$, restricts to properties that involve neighbors of a given element. 
These logics are most naturally applied to finite models, to avoid dealing with addition or parity of infinity: 
thus in this work, \emph{finite satisfiability will always be our default}, and 
we just refer to ``satisfiability'' henceforward. 
In~\cite{usicalp} it was shown that local Presburger logics are in some sense equally expressive as 
certain variations of Graph Neural Networks (GNNs), and 
GNN verification problems can be reduced to satisfiability problems for these logics.
It has been shown that satisfiability of Presburger logic is undecidable, while satisfiability of local Presburger logic is decidable~\cite{percentage}: but the exact complexity of local Presburger logic had not been determined. Another way of enhancing logic with arithmetic involves \emph{global cardinality constraints}. 
In such an enhancement, we can write a constraint that defines a constraint (e.g., a linear inequality) 
between one variable formulas. But we cannot have such a constraint parameterized by a second free variable.  For example, we can write that unary predicate $U$ has cardinality below unary predicate $V$. But if $E$ and $F$ are two binary predicates, we cannot express that the elements that are not $E$ adjacent to $x$ have greater cardinality than the elements that are not $F$ adjacent to $x$.

As with local Presburger logic, it follows from prior results  that $\Ctwo$ enhanced with global cardinality constraints has decidable (finite) satisfiability problem.
The exact complexity is not explicitly studied in the literature,  although in the case of global cardinality constraints consisting of linear inequalities, tight bounds follow from results of Rudolph in~\cite{rudolphcardinality}.    A natural approach to proving decidability for logics
enhanced with global constraints follows is to show that the \emph{the spectrum of each tuple of unary formulas of the logic is effectively semilinear}. The spectrum of a tuple of unary formulas $\phi_1(x) \ldots \phi_k(x)$ is the set of possible values that the tuple of cardinalities of the satisfiers of the formula can have within a finite model. For example, for the pair of formulas $\phi_1(x)=U(x), \phi_2(x)= U(x) \wedge V(x)$ the spectrum consists of all pairs $(m, n)$ with $n \leq m$. The spectrum being effectively semilinear means that we can find a quantifier-free formula of arithmetic that defines it. From effective semilinearity of the spectrum, we easily obtain decidability with additional global constraints.
Here we provide  complexity bounds for several of these problems, one featuring  arithmetic over counts of neighbors, another concerning computing the spectrum of formulas. 
See Table~\ref{tab:results} for the summary of our new upper bounds.
At the same time, we provide simpler proofs of two fundamental results in the area: decidability of (finite) satisfiability for $\GPtwo$ (see Section~\ref{sec:gptwo}), and effective semi-linearity of the spectrum for $\Ctwo$ (Section~\ref{sec:ctwo}). We complement these positive results with a counterexample to semilinearity of spectra for $\GPtwo$ unary formulas (Section~\ref{sec:gptwocounter}).
